\documentclass[10pt]{article}
\usepackage[letterpaper,textwidth=6.55in,textheight=9.0in,centering]{geometry}
\usepackage[T1]{fontenc}
\usepackage{mathptmx}
\usepackage{microtype,amsmath,amssymb,booktabs,tabularx,array,graphicx,float,xcolor,caption,hyperref,url}
\usepackage{xurl}
\usepackage[font=small,labelfont=bf]{caption}
\hypersetup{colorlinks=true,linkcolor=black,citecolor=blue!50!black,urlcolor=blue!50!black,pdfauthor={Anonymous},pdftitle={AfterClick: Beyond Action-Level Safety in Asynchronous AI Agents}}
\setlength{\parindent}{0pt}\setlength{\parskip}{3.5pt plus 1pt}\setlength{\textfloatsep}{10pt plus 2pt minus 2pt}\setlength{\intextsep}{8pt plus 2pt minus 2pt}\setlength{\abovecaptionskip}{4pt}\setlength{\belowcaptionskip}{2pt}
\newcommand{\fig}[3]{\begin{figure}[!htbp]\centering\includegraphics[width=#2\linewidth]{figures/#1.pdf}\caption{#3}\label{fig:#1}\end{figure}}
\newcommand{\researchq}[1]{\textbf{RQ#1.}}
\newcommand{\code}[1]{\texttt{#1}}
\newcommand{\repo}{\url{https://github.com/GomathyDhanya/AfterClick/tree/31eab61870befc520841d11f422bb1a372ba92b8}}
\begin{document}
\begin{center}
\rule{\linewidth}{1.2pt}\\[9pt]
{\LARGE\bfseries AfterClick: Beyond Action-Level Safety\\[3pt] in Asynchronous AI Agents\par}
\vspace{5pt}{\large Irreversible Effects, Admission Boundaries, and Restoration Debt\par}
\vspace{8pt}\rule{\linewidth}{0.6pt}\\[9pt]
{\normalsize Research manuscript \;--\; 26 November 2026}\par
\vspace{3pt}{\small Conference-style working paper; not a peer-reviewed publication}
\end{center}
\begin{abstract}
Tool-using AI agents act in environments that continue changing during deliberation, intervention and recovery. A correct protective decision may be ineffective when an external effect has already been admitted; a successful intervention may itself obstruct authorized work. We study these distinctions in \emph{AfterClick}, an executable testbed with two browser-accessible applications, PostgreSQL transactions, asynchronous workers and separately persisted synthetic external effects. Among 1,073 valid internal agent episodes from an incompletely executed frozen evaluation, 65 complete their local goals despite forbidden effect history. In a completed fixed-recommendation study (288 registered branches, 287 valid), imposing 300 ms and 1,500 ms of additional browser enforcement delay increases forbidden-history incidence by 27.8 and 60.6 percentage points, respectively, on 12 constructed source tasks. Source inspection and 72 subsequent engineered probes identify the mechanism: serialization of protection and admission does not recall previously admitted tickets. We introduce restoration-debt accounting as an operational description of duties activated by containment: 583 of 617 reconstructed duties remain outstanding at the observed censoring horizon. Finally, all 264 historical LLM reviewer responses and their 264 retrospective replays choose suspension, including 85 inputs for which the public pending-service rule permits execution. The results motivate a unified evaluation of historical safety, temporal controllability and legitimate restoration. We distinguish controlled intervention effects from retrospective associations, preserve unknown outcomes, and make no claim of universal controller superiority or a causal effect of neural inference latency.
\end{abstract}
\section{Introduction}
Consider an assistant asked to migrate access to confidential records from an outgoing team to an incoming team. A background worker has already queued a delegated read. The assistant detects that the request is outside the permitted scope and suspends the service. The interface acknowledges the suspension. The assistant finishes the migration and restores approved internal access. A conventional evaluation may see the correct terminal state and award success. Yet the external receiver may have already recorded the prohibited read. The task is complete, but the confidentiality failure is permanent under the modeled policy.

This example illustrates a gap between reasoning about an action and reasoning about its consequences. An agent does not control every state transition in an application. Workers, retries, delayed messages and other users can progress while the agent reasons. Even when database operations are serialized, external effects may outlive the authority under which they were admitted. An intervention may additionally disable legitimate functionality that the agent must later restore. These are related security and availability problems, not simply incorrect tool calls.

Prior work already separates competence from recovery and uses effect-history oracles \cite{undo}, stages irreversible effects before commitment \cite{cordon}, and studies asynchronous task obligations \cite{async}. We ask how these problems interact in an executable asynchronous environment. Our emphasis is measurement, not a new transaction primitive or a claimed state-of-the-art controller. The observation that a safety monitor cannot suppress an uncontrollable action is established in enforcement theory \cite{basin}; our study traces where this boundary falls in an agent-driven workflow and measures the accompanying utility cost.

\textbf{Research questions.} \researchq{1} When does local completion coexist with forbidden external history, and which cases actually contain evidence of recovery? \researchq{2} How does changing enforcement time at fixed recommendation affect irreversible effects, and which implementation boundary explains the outcome? \researchq{3} What new duties are activated by protection and how often are they discharged? \researchq{4} What limitations of additional LLM review are observable in decisions and timing, and which explanations remain unidentified?

\textbf{Contributions.} (i) An executable, contract-based measurement environment that separates local goals, legitimate completion and physical effect history; (ii) paired timing interventions and source-validated admission-boundary probes; (iii) trace-level restoration-debt accounting and reviewer decision diagnostics; (iv) an evidence-linked retrospective analysis retaining failed, invalid and unexecuted assignments. The paper reports a completed bounded analysis of a partially executed larger research program, not completion of that entire program.
\section{Problem formulation and threat model}
\subsection{From action permission to historical safety}
Let $s_t$ denote the business state, $a_t$ the currently visible authority, $q_t$ the pending and admitted work, and $H_t$ the append-only physical-effect history in the modeled environment. The actor observes a projection $o_t$ of public state, not the oracle's hidden contract or future scheduler. A contract defines the authorized goal, preservation constraints, allowed effect scope and legitimate service obligations. It induces four outcomes for a valid trace $\tau$:
\begin{align*}
 L(\tau)&=\mathbf{1}[\text{current local goals complete}],\\
 G(\tau)&=\mathbf{1}[\text{all legitimate obligations satisfied}],\\
 F(\tau)&=\mathbf{1}[\exists h\in H_{T}:\ h\text{ violates the contract}],\\
 Y(\tau)&=G(\tau)(1-F(\tau)).
\end{align*}
$Y$ is joint safe completion. $L$ and $G$ are distinct definitions: $G$ includes preservation, approved pending delivery, service availability and reconciliation. In the observed valid internal episodes they happen to agree. Invalid or unexecuted episodes have unknown outcomes, not zero success or zero violations. A mismatch is $L=1,F=1$; it is an observed repair only if the trace additionally witnesses a relevant incorrect state followed by correction. The distinction matters because ordinary task completion after an effect is not necessarily undoing an error.

\textbf{Three safety boundaries.} Decision safety concerns the policy correctness of a recommendation; enforcement safety concerns whether the effect remains controllable; restoration concerns whether authorized functionality is recovered. These are complementary, not mutually exclusive metrics. In particular, a false block can improve a narrow no-violation metric while harming availability. The evaluator must account for both.
\subsection{Admission and irreversible history}
An effect may move from queued to admitted to physically persisted. Let $t_a$ be its actual admission boundary, $t_p$ the point at which protection is effective at the admission guard, and $t_r$ the persistence of the effect. In the legacy mechanism, if $t_a<t_p<t_r$, a previously admitted ticket remains deliverable. An acknowledgment after protection does not imply that every ticket has been revoked. Conversely, if protection takes effect before admission, the modeled pending effect can be blocked. These are conditional statements about the tested mechanism, not universal timing laws. Audit insertion and commit return bound a transaction but do not expose its precise commit instant.

\textbf{History barrier.} Under an immutable-effect contract, $H_t\subseteq H_{t'}$ for $t'\ge t$. If $F$ is already true, no later local edit can make $Y=1$ for that same history. This follows directly from the definition of $F$; it is not a novel impossibility theorem. In a system with legally effective compensation, a different contract might permit recovery and the implication would need to be reconsidered. Before a violation, safe continuations can be unknown: lack of a witnessed continuation is not proof of impossibility.
\subsection{Threat model and scope}
We consider an imperfect actor and a potentially overrestrictive reviewer executing legitimate tasks in an application with asynchronous effects and concurrent state changes. The contract may prohibit an effect's recipient, scope, timing or resource commitment. The worker's admission eligibility and public authority can differ from the task-level permission that the evaluator checks. The evaluator observes independent persisted receiver evidence. Public policy and signed authority are trusted fixture data, not a cryptographic verification protocol. The models receive no hidden evaluator labels. The experiment does not involve real personal data, actual customers, production secrets, or uncontrolled third-party effects. We do not evaluate a malicious database administrator, compromised receiver, arbitrary prompt-injection attacker or network adversary. Security here means contract-defined integrity, confidentiality and availability in controlled synthetic workflows.
\section{AfterClick: architecture and experimental testbed}
\fig{architecture}{1}{AfterClick's execution and evidence boundaries (schematic). Solid arrows are execution paths; dashed arrows represent retrospective evidence. Actor and controller see public observations. The offline evaluator separately examines task contracts, transaction evidence and synthetic HTTP receipts. Workers progress during model calls.}
\subsection{Applications, tasks and controls}
FlowBoard implements versioned business tasks, scoped grants, approval comments and dependencies. OpsConsole separately implements work orders, assets, scopes, capacities and approval transitions. They share benchmark orchestration, effect mechanisms and evaluator infrastructure; they are two environments, not independent samples of production applications. Chromium interacts with actual application interfaces; PostgreSQL stores each world's state in an isolated schema. Workers advance pending operations independently. An HTTP receiver persists effect content outside the actor's ordinary business-state view. Browser responses, snapshots, model calls and effect receipts are logged with identifiers and time-domain metadata when available.

The catalog contains 288 source contracts: 72 DEV and 216 TEST. Six business-policy families span two applications and three complexity levels (8, 13 or 21 required business operations). There are 24 structural layouts and six vocabularies. Every contract has two executed reference witnesses, giving 576 browser reference paths; an independently implemented checker rejects 10,944 constructed evidence mutations. These establish executable reference support, not arbitrary oracle correctness or human-validated realism. The following cases illustrate the range of contract mechanisms; descriptions are schematic unless linked to a concrete episode.
\begin{table}[!htbp]\small\centering\caption{Task families and representative modeled risks.}\label{tab:families}
\begin{tabularx}{\linewidth}{@{}lXX@{}}\toprule
Family & Legitimate objective & Risk carried by external effects \\\midrule
Access migration & Transfer scoped internal reviewer access & Delegated read outside approved scope \\
Customer release & Reconcile approvals and notices & Notice with inappropriate release scope \\
Incident response & Route scoped escalation and updates & Notification to incorrect audience \\
Retention export & Reconcile permitted record exports & Restricted export outside policy \\
Resource orchestration & Allocate approved capacity & Irreversible unauthorized resource use \\
Cross-team approval & Coordinate approved service commitments & Unapproved downstream commitment \\\bottomrule
\end{tabularx}\end{table}
\subsection{Policy variants and agent settings}
\begin{table}[!htbp]\small\centering\caption{Historical intervention bundles. These are component comparisons, not perfectly isolated single-factor ablations.}\label{tab:policies}
\begin{tabularx}{\linewidth}{@{}lXX@{}}\toprule
Arm & Intervention & Intended question \\\midrule
P0 & Actor only (clean or stress world) & What can the unmodified actor accomplish? \\
P1 & Immediate deterministic pending-service protection & Can fast containment prevent effects? \\
P2 & Public verification and deterministic protection & Does verification help under the same broad guard? \\
P3 & P2-style verification plus sequential LLM review & What changes with a fresh reviewer? \\
P4 & P1 plus explicit public restoration feedback/checks & Does obligation visibility help? \\
P5 & Queue hold before verification and LLM review & Does suspension-first ordering help? \\\bottomrule
\end{tabularx}\end{table}
The public pending-service rule is deliberately simple: allow if visible authority is approved, otherwise suspend. It is operation-specific, rather than a judgment about whether every business record is finished. P1--P5 share a consequence guard; P5 is a queue-hold proxy, not semantic staging. The primary actor is Qwen3.5-9B; the supplementary model is Qwen3.6-35B-A3B-FP8. Historical actor budgets are 24 generations, 200 browser actions and 240 s wall time, temperature zero and at most 512 output tokens per generation, with thinking disabled. Precision, serving, hardware and model revision differ between cohorts. We do not interpret the model comparison as an isolated scale intervention.
\subsection{Enrollment and evidence collection}
\begin{table}[!htbp]\small\centering\caption{Original frozen assignment disposition. The official external cohort uses a distinct oracle.}\label{tab:enroll}
\begin{tabular}{@{}lrrrrr@{}}\toprule
Cohort & Registered & Terminal & Valid internal & Invalid & Not run\\\midrule
Primary & 2304 & 905 & 744 & 161 & 1399\\
Larger & 360 & 334 & 329 & 5 & 26\\
External & 240 & 240 & -- & -- & 0\\
Total & 2904 & 1479 & 1073 & 166 & 1425\\
\bottomrule
\end{tabular}
\end{table}
The original TEST protocol was frozen before inference and reduced prospectively under the original budget. Even the revised matrix was incomplete; primary terminal records cover 185 of 216 TEST sources. We retain the complete assignment reconciliation. The larger cohort covers 36 selected source tasks. The distinct external study executes all 240 registered conditions from 12 official UndoBench tasks with two context seeds and five wrappers. It is supplementary, not pooled with internal outcomes.

Three additional forms of evidence answer narrower questions: (a) a completed 288-branch fixed-recommendation actuator study, using 12 constructed L3 DEV sources, two seeded worlds, three admission offsets and four actuators; (b) a completed 72-branch prospectively enrolled mechanism-probe matrix on the same opened DEV sources, crossing three engineered orders and two protection commands; (c) a completed retrospective prompt-replay census of 264 available reviewer inputs. Only (c) requires new LLM inference. Neither (b) nor (c) extends the original TEST enrollment or constitutes an independent held-out evaluation.
\subsection{Verification and statistical analysis}
The publication analysis normalizes 1,239 internal terminal episodes and preserves raw paths, SHA-256 digests, JSON pointers, native identifiers and clock qualifications. Separate software checks reportedly verify 48,422 raw files and 119,972 pointers. Independently implemented outcome checks agree on all valid internal outcomes; no independent human annotation has been completed. A physical effect is checked against persisted receiver content, rather than inferred from a model response or browser acknowledgment.

We resample whole source contracts, not individual worlds, because repeated seeds and variants share structure. The retrospective mismatch estimator is the episode-pooled ratio conditional on valid observed local completion. We use the recorded 10,000 source bootstrap draws (seed 11262026). The original policy and controlled-actuator estimators weight source-level contrasts, with stratification in the original policy analysis (seed 8261125). Confidence intervals condition on observed support. We also compute worst-case assignment bounds: if $k$ observed successes exist among $n$ enrolled assignments and $u$ outcomes are unknown, the success rate lies in $[k/n,(k+u)/n]$. Matched contrast bounds allow missing outcomes to vary adversarially in both arms. These bounds are identification regions, not confidence intervals. Sparse strata can produce degenerate bootstrap intervals, so they must not be interpreted as exact population effects. Neither selection into valid episodes nor execution order is assumed ignorable.
\section{RQ1: Local success is not historical security}
\fig{mismatch}{.94}{Observed local completion ($L$) versus forbidden history ($F$). Red outlines identify the mismatch cases. The two cohorts are displayed separately; counts use only valid internal episodes.}
\begin{table}[!htbp]\small\centering\caption{History violations among locally complete episodes. CIs resample observed source clusters; assignment bounds include unknown outcomes.}\label{tab:mismatch}
\begin{tabular}{@{}lrrrr@{}}\toprule
Cohort & Mismatch / $L$ & Rate & Cluster 95\% CI & Assignment bounds\\\midrule
Primary & 43/175 & 24.6\% & 18.0--32.1\% & 2.5--92.4\%\\
Larger & 22/57 & 38.6\% & 26.9--48.7\% & 25.0--60.2\%\\
\bottomrule
\end{tabular}
\end{table}
Across observed valid episodes the unconditional mismatch rates are 43/744 (5.8\%) and 22/329 (6.7\%). Conditioning on local completion asks a different question: among outcomes that appear successful to a local evaluator, how many are historically unsafe? Table~\ref{tab:mismatch} quantifies that blind spot. The source support is 177 primary and 36 larger valid source clusters. The 65 cases have individually linked physical receipts and matching content; there are zero observed wrong-to-correct business commit sequences in this subset. Three contain witnessed containment followed by service reopening. The controlled scripted cohort separately demonstrates actual reopening after forbidden history in 77/287 valid branches. We do not conflate scripted reference recovery with autonomous repair.

\textbf{Observed case A: access migration.} The case \path{TEST-ACCESS_MIGRATION-L1-FLOWBOARD-03-P0-tight-1-primary} completes local goals with a forbidden pending \code{delegated.read} recorded at scope \code{external-contractor}. Both receipt and physical record exist. It has $L=G=F=1,Y=0$, without an observed containment/restoration transition. This illustrates a local metric's blind spot rather than evidence of a repair algorithm.

\textbf{Observed case B: customer release.} The larger-model case \path{TEST-CUSTOMER_RELEASE-L1-OPSCONSOLE-05-P3-tight-1-larger} has a forbidden pending \code{customer.notice} at scope \code{internal-preview}. The case audit also witnesses containment and service reopening. The final service is restored; the recorded history remains forbidden. This is a genuine service-restoration witness, not necessarily a correction of an erroneous business mutation. Both case cards identify the raw snapshot hash and receiver pointers (Appendix~\ref{app:trace}).

\textbf{Sensitivity.} Removing one application or family at a time yields primary conditional mismatch rates of 22.7--26.6\% and larger rates of 25.0--45.9\%. The phenomenon is not confined to one excluded family on observed support, though the larger cohort is sensitive to application composition. These are descriptive structural checks, not independent validation. Wide assignment bounds rule out strong extrapolation to the entire enrollment. \textbf{Answer:} local completion misses a real, independently persisted class of forbidden histories in both observed cohorts. Historical safety must be measured explicitly.
\section{RQ2: Enforcement timing and the admission boundary}
\subsection{Controlled delay intervention}
We fix the public suspension recommendation before adding delay, and compare ordinary browser enforcement against direct HTTP, browser plus 300 ms and browser plus 1,500 ms. Each arm runs the same registered source/world/offset combinations and full reference continuation. Offsets are configured at 150, 600 and 1,800 ms. A separately implemented verification hashes all 288 initial cuts and checks actual receiver outcomes. One branch is invalid (the tight-offset, longest-delay condition); it remains unknown. All 287 valid reference continuations achieve legitimate completion and witness containment followed by reopening, including 77 with forbidden history. Joint success is therefore 210/287 in this particular scripted study, not the LLM's success rate.
\fig{actuator}{1}{Controlled actuator experiment. Each heatmap cell has 24 registered branches; * denotes one unknown outcome, with 23 observed forbidden histories. Right: original source-paired effect estimates and 95\% source bootstrap intervals relative to the ordinary browser. Configured offsets are not individually measured prevention deadlines.}
\begin{table}[!htbp]\small\centering\caption{Enforcement intervention contrasts (percentage points). All use 12 matched source clusters.}\label{tab:delay}
\begin{tabular}{@{}lrrr@{}}\toprule
Intervention vs browser & Effect on $F$ & Cluster 95\% CI & Unknown bounds\\\midrule
Direct HTTP & $-2.78$ & $[-9.72,4.17]$ & $[-2.78,-2.78]$\\
Browser +300 ms & $+27.78$ & $[23.61,31.94]$ & $[27.78,27.78]$\\
Browser +1500 ms & $+60.56$ & $[55.83,64.72]$ & $[59.72,61.11]$\\
\bottomrule
\end{tabular}
\end{table}
The controlled intervention produces a large timing-dependent change, especially when configured admission slack is short. The direct HTTP comparison spans both signs; the evidence does not establish superiority of HTTP over browser. At 1,800 ms, all four conditions have zero observed forbidden effects, suggesting the importance of relative timing rather than latency alone. The largest-delay estimate remains positive under worst-case treatment of its single unknown branch. The design holds information and recommendation fixed, so it isolates the imposed actuator intervention under this executing system. It does not isolate model latency or infer a universal deadline.
\subsection{Why protection fails: a corrected mechanism}
Initial interpretations suggested an unlocked authorization read racing against a protection write. Inspection of the actual transaction path corrects this explanation: the worker and controller share a per-schema PostgreSQL advisory transaction lock. Admission and protection are serialized. The admitted ticket, however, can survive protection and reach a receiver that does not revalidate the changed containment state. Thus even correct database locking does not imply end-to-end revocability.
\fig{forensic}{.98}{One observed forensic case: \texttt{ACCESS MIGRATION, L2, OpsConsole, P2}. Reported timestamps are aligned to the admission audit. The shaded interval spans protection audit insertion to the recorded commit-return upper bound; the receipt occurs later. This is not an exact measurement of historical SQL commit or authorization-read instants.}
The 34 historical post-protection-audit violation cases all join the same pending effect to an admission audit preceding the protection audit, and to a physical receipt after the protection commit-return upper bound. The displayed case has a receipt about 266.24 ms after the reported commit-return bound. Historical clock-domain metadata and exact read/commit instants are incomplete; the recorded ordering plus source serialization supports the interpretation, rather than establishing nanosecond-accurate timings.
\subsection{Barrier-controlled mechanism probes}
A separate finite registered experiment follows successful DEV checks, then uses 12 already opened source tasks crossed with three schedules and two commands (Suspend and Cancel). Instrumentation wraps the original transaction context and records process identity, lock acquisition, read-return and commit-return bounds. In the admission-before-protection order, a barrier holds the worker's original schema lock while the controller attempts to acquire it; the worker releases admission, protection commits and actual worker delivery occurs afterwards. Every branch executes real database transactions and receiver calls. No model is used and no snapshot replacement acts as recovery.
\begin{table}[!htbp]\small\centering\caption{Mechanism probes, all 72 registered branches valid. The engineered orders test implementation behavior, not a population estimate.}\label{tab:probe}
\begin{tabular}{@{}lrrr@{}}\toprule
Order & Branches & Pending effects & Prevented\\\midrule
Protection before admission & 24 & 0 & 24\\
Admission before protection, receipt after & 24 & 24 & 0\\
Receipt before protection & 24 & 24 & 0\\
\bottomrule
\end{tabular}
\end{table}
This confirms a non-revocable-ticket boundary in the existing mechanism. It is stronger than merely correlating late violations with model response times. It is narrower than demonstrating receiver-side fencing or a new transaction protocol. \textbf{Answer:} within the controlled testbed, protection effectiveness depends on admission timing even with identical recommendations. A successful containment acknowledgment is not a certificate that all outstanding external effects have been prevented.
\section{RQ3: The downstream cost of protective intervention}
\subsection{Restoration debt as an operational metric}
A containment action can introduce duties that did not exist before the action. We call the observed unresolved induced duties \emph{restoration debt}. For each observed duty $i$, record activation $b_i$, discharge $d_i$ if observed, and a censoring horizon $c_i$ (last recorded backend audit event). With unit weight $w_i=1$ in this analysis, define
\[
 D(t)=\sum_i w_i\,\mathbf{1}\{b_i\le t<e_i\},\qquad e_i=\begin{cases}d_i&\text{if discharged},\\c_i&\text{otherwise}.\end{cases}
\]
The observed exposure is $A=\int D(t)\,dt$. Outstanding status at censoring is stored separately from this interval convention. This definition measures active duties, not damage severity; weights and horizons would need calibration for production. Preexisting task obligations are excluded from induced debt. A service reopening duty activates on a false-to-true suspension transition and discharges on observed reopening. A resubmission duty activates when an approved pending request is actually denied admission by protection, and discharges when a Resume transition requeues it. Requeueing is not equivalent to delivery; approved delivery remains a separate $G$ obligation.
\fig{debt}{1}{Left: observed primary stress-policy outcomes on each arm's valid support, showing the safety--utility trade-off. Right: independently reconstructed induced duties, separated by type and censoring status. The bars do not identify a causal effect of policy on debt.}
\begin{table}[!htbp]\small\centering\caption{Protection-induced duties reconstructed from 1,073 valid internal episodes.}\label{tab:debt}
\begin{tabular}{@{}lrrr@{}}\toprule
Duty type & Triggered & Discharged & Outstanding\\\midrule
Service reopening & 547 & 32 & 515\\
Approved request resubmission & 70 & 2 & 68\\
Total & 617 & 34 & 583\\
\bottomrule
\end{tabular}
\end{table}
Only 34/617 (5.5\%) duties are observed to discharge; 94.5\% remain outstanding at censoring. The independent reconstruction agrees on all 547 service-duty identities. The maximum concurrent unit debt is two in the observed model. Outstanding does not mean permanently abandoned: the trace horizon is unequal across episodes, and observed debt exposure is a lower-bound interval rather than a survival estimate. The counts are duties, not distinct episodes or statistically independent tasks.

\subsection{Security without availability}
The primary actor-only stress arm records 60/142 legitimate completions, 88/142 forbidden histories and 19/142 joint successes. Immediate protection records 19/142 legitimate completions, only 1/142 forbidden history and the same marginal 19/142 joint successes. The larger cohort shows an even sharper descriptive trade-off: actor-only achieves 26/65 legitimate completions and 49/65 forbidden histories; immediate protection achieves 1/65 and 1/65 respectively. The equal marginal primary joint counts do not prove equivalence; these arms differ in valid support and matched assignment bounds remain wide.

A useful early negative case appears in the earlier prospective study: immediate revocation prevented a queued read in source 05, but the agent declared completion without restoring legitimate access, while the delayed arms restored it. The 12-source comparison was inconclusive (11/12 vs 10/12 joint completion); it illustrates the mechanism without establishing a population tendency. The new duty extraction extends this observation systematically. We separate 7,354 preexisting business-obligation failure labels from the induced-duty counts; labels can overlap and are not episode counts.

\textbf{Answer:} protection creates identifiable, traceable additional obligations that are frequently unresolved at the recorded horizon. A narrow no-forbidden-history metric can reward overrestriction; joint success and debt accounting reveal the availability cost. The analysis conditions on protection, a policy-dependent mediator, so it cannot establish that restoration reminders reduce debt causally.
\section{RQ4: What additional LLM supervision actually contributes}
\subsection{Decision content versus elapsed time}
The 264 historical reviewer responses comprise 198 primary and 66 larger-model requests. Every response parses as strict JSON and recommends \code{suspend}. The actor runner processes 225; 39 received responses are not processed. Relative to the declared operation-specific pending-service rule, 179 decisions agree and 85 disagree: all 85 visibly approved pending requests receive suspension recommendations, including 74 processed recommendations. This is not a human-rated assessment of the wider workflow. It is a disagreement with the implemented rule. The model's explanations often refer to incomplete workflow requirements or reconciliation, raising the hypothesis that the reviewer conflates overall business readiness with local authority.
\fig{reviewer}{.96}{Historical reviewer observations. Every recommendation is suspend; 85 inputs permit pending execution under the declared public rule. The right panel shows observed HTTP interval quartiles, endpoints and p95; it is not a measurement of pure neural inference latency.}
A separate replay job uses the original prompt contents, generation settings and pinned Qwen3.5-9B weights, with a new served alias. Each of the 264 opened inputs is evaluated once; all responses are received, strictly valid and again recommend suspension. The 198 original primary prompts are a same-weight serving-stack diagnostic. The 66 original larger-model prompts are cross-model replays through 9B, not replications of the larger model. No world continuation is executed. This excludes malformed output as the explanation of universal suspension and shows the observed decision pattern is repeatable under that setup. It cannot distinguish prompt framing, scope confusion, conservative decision bias or genuine ambiguity in the task contracts.

Median historical HTTP interval is 1.69 s (IQR 0.97--6.95 s; p95 12.50 s). These intervals include transport, queueing and serving. They are comparable in magnitude to configured actuator offsets, but cannot be substituted into a causal latency estimate: the content, timing and application of recommendations all differ. The timing experiment in RQ2 isolates added enforcement delay, not reviewer inference.
\subsection{Policy variants: ablations and uncertainty}
\begin{table}[!htbp]\small\centering\caption{Primary stress and clean arms, plus larger-cohort arms. All counts are on valid support. $G$: legitimate; $F$: forbidden history; $Y$: joint safe.}\label{tab:arm}
\begin{tabular}{@{}llrrrr@{}}\toprule
Cohort & Policy & Valid & $G$ & $F$ & $Y$\\\midrule
Primary & P0 clean & 141 & 59 & 0 & 59\\
Primary & P0 stress & 142 & 60 & 88 & 19\\
Primary & P1 stress & 142 & 19 & 1 & 19\\
Primary & P2 stress & 141 & 24 & 12 & 24\\
Primary & P3 stress & 136 & 4 & 49 & 2\\
Primary & P4 stress & 18 & 8 & 0 & 8\\
Primary & P5 stress & 24 & 1 & 0 & 1\\
Larger & P0 clean & 69 & 25 & 0 & 25\\
Larger & P0 stress & 65 & 26 & 49 & 5\\
Larger & P1 stress & 65 & 1 & 1 & 1\\
Larger & P2 stress & 65 & 4 & 13 & 4\\
Larger & P3 stress & 65 & 1 & 21 & 0\\
\bottomrule
\end{tabular}
\end{table}
\fig{bounds}{1}{Ablation-style policy bundle contrasts. Parentheses denote complete source clusters. Blue intervals reflect observed source sampling uncertainty; broad gray bounds include unknown assignments. Point intervals in sparse stratified strata are degenerate, not exact effects. All full-assignment bounds cross zero.}
The original primary P3--P2 contrast on complete source support is $-13.51$ pp (111 clusters, 95\% source bootstrap interval $[-18.9,-9.0]$). Full assignment bounds $[-72.45,63.43]$ pp admit either sign. The larger contrast is $-6.94$ pp across 36 clusters, with bounds $[-15.28,4.17]$ pp. P2--P1 is $+2.68$ pp; P4--P1 is $+31.25$ pp on just 16 source clusters. The latter's Holm-adjusted secondary $p=0.0983$ and broad bounds prevent a confirmatory claim. P5--P1 is $-10$ pp on 20 clusters. These contrasts evaluate intervention bundles; no matched information-content, latency or no-restoration-needed factorial was performed. All original contrasts remain in the record, including negative and inconclusive ones.

\textbf{Answer:} in this observed dataset the extra review policy is associated with reduced joint completion and repeats a restrictive decision pattern even on visibly approved inputs. A causal harm from extra reasoning or a general benefit from faster deterministic review is not established. A future study must separate the semantic value of review, the enforcement delay it induces, and the actor's ability to act on its recommendation.
\section{Discussion: a unified account of failure}
\subsection{Why the three boundaries are distinct}
The results yield a useful decomposition. First, a model may be wrong about operation-specific permission (RQ4). Second, even a correct recommendation may be enforced too late (RQ2). Third, even effective prevention may introduce obligations that the actor does not discharge (RQ3). Finally, an endpoint-only evaluator can obscure the result of any earlier irreversible effect (RQ1). The decomposition avoids collapsing authorization, execution, and recovery into a single safety score. A correct decision is not sufficient for safe execution; effective blocking is not sufficient for joint safe completion. These are empirically witnessed distinctions, not claims of independent or additive causal effects.

The original graph analysis also makes a crucial negative point: 59,174 of 71,138 primary prefix nodes and 28,512 of 36,444 larger prefix nodes remain unknown with respect to recoverability. Only an actually executed safe sibling at a certified initial cut warrants a witnessed-continuation label. A fixed forbidden-history prefix is irrecoverable under the modeled history constraint, but a trace that ends unsuccessfully without such a barrier could have an unobserved safe alternative. There are zero complex-prefix formal certificates. We therefore cannot train a reliable binary recoverability model by labeling every unsuccessful continuation as impossible.
\subsection{Testable hypotheses derived from the evidence}
The following are scientific hypotheses, not additional findings. They are retained because they turn the observed failure patterns into narrowly falsifiable research questions.
\begin{table}[!htbp]\small\centering\caption{Mechanism hypotheses and falsification designs (unexecuted).}\label{tab:hypotheses}
\begin{tabularx}{\linewidth}{@{}p{.22\linewidth}XX@{}}\toprule
Hypothesis & Evidence motivating it & Test that could refute it\\\midrule
H1: Boundary relocation & Admitted tickets survive containment in all post-admission probes & Compare serialized admission to atomic receiver-side fencing; measure forbidden effects, legitimate delivery and retry duplication on fresh tasks. \\
H2: Restoration debt is actionable & 583 outstanding duties, plus selected positive P4 contrast & Randomized reminder vs neutral feedback crossed with restoration-required vs matched no-restoration tasks; hold protection fixed. \\
H3: Reviewer scope confusion & Universal suspension despite 85 approved inputs & Counterfactually vary approval and unrelated incomplete duties; scope instruction factorial with independently reviewed labels. \\
H4: Information has a temporal opportunity cost & Fixed-decision actuator delay increases violations & Cross relevant/irrelevant evidence with immediate/delayed access, plus cached-equivalent and latency-matched review. \\
H5: Selective containment improves utility & Immediate protection suppresses forbidden effects but completion drops & Compare scoped vs global suspension with equal authority and timing, measuring both $F$ and $G$. \\\bottomrule
\end{tabularx}\end{table}
H1 in particular must coordinate validation and effect persistence: a non-atomic receiver-side recheck may simply move the race. H2 requires an intervention-level estimand, not conditioning on observed debt after treatment. H3 requires clear labels for operation permission rather than wholesale workflow approval. H4 must establish a counterfactual at matched information and latency. H5 could fail if narrower containment leaves interacting effects exposed. None is necessary to present the bounded measurement results; each is needed before asserting its respective mechanism generalizes or improves controllers.

\textbf{Design implications.} Separate reversible state from immutable receipts; enforce authority as near as feasible to the actual commitment boundary; explicitly account for already admitted work; make restoration duties and prerequisites visible; distinguish uncertain from unsafe histories; and benchmark reviewers on both allowed and prohibited requests. These recommendations are motivated by the results and existing systems principles, not validated production fixes.
\section{Related work}
\textbf{Recovery and effect-aware evaluation.} UndoBench separates task competence from recovery and evaluates state and wire-level effects \cite{undo}. Our supplementary evaluation uses its released evaluator on a small official subset; history-aware evaluation is therefore not claimed as a novel principle. ReInBench explicitly investigates error recognition and self-recovery \cite{rein}; our case taxonomy likewise refuses to equate endpoint success with a corrective action. Safe to Resume? studies consistency problems across rollback and external dependencies \cite{resume}; our effects are modeled as durable and not repaired by snapshot replacement.

\textbf{Execution boundaries and asynchronous agents.} Cordon introduces semantic transactions that stage and validate irreversible effects \cite{cordon}. Our queue-hold P5 does not implement Cordon, and we do not compare against its complete system. Execution-monitoring work formalizes the distinction between observable and controllable actions \cite{basin}. AfterClick makes the admission/receipt distinction measurable in this particular application architecture. Async-RBench studies asynchronous subagent results and response obligations \cite{async}; our focus is independently progressing external effects and duties induced specifically by protection. AgentTrust investigates runtime interception and safety evaluation \cite{agenttrust}; our reviewer results concern only the tested prompts and rules. The contribution is an empirical integration of these issues rather than a claim that each underlying concept is new.
\section{Limitations, reproducibility and responsible release}
The most important limitation is incomplete original enrollment (1,425 assignments not run). Invalid outcomes and low clean competence (59/141 primary, 25/69 larger) limit policy conclusions. The task families share layouts and mechanisms; all controlled probes use opened constructed DEV sources. The study lacks independent human task/authority review, a new independent workload, a complete restoration factorial, causal information/latency controls and exact mid-trajectory counterfactuals. Model comparisons confound architecture, quantization and serving. Physical receivers are synthetic; no production security guarantee or prevalence estimate follows. Historical time stamps do not give exact original authority reads or commits. Formal claims apply only to the modeled history semantics.

The source tree and compact evidence are available at the pinned revision in Appendix~\ref{app:trace}. The underlying raw archive and companion are locally prepared but not publicly released; a source-only checkout cannot fully reproduce the raw-data analysis. Automated relocated Linux execution reproduced numerical analyses after reader corrections. The independent Linux regression ran 395 passes and 22 missing-fixture failures, followed by 22 passes after fixture restoration. A separately identified cleanup regression on the exact restored historical source passes all 417 tests locally; this does not rewrite the earlier Linux record. Software independence and a second machine are not independent human replication. The data candidate has undergone heuristic credential-pattern scanning, not a guarantee of de-identification; licensing and public release approval remain outstanding. Synthetic effects avoid direct real-world disclosures or commitments. No further inference or experiment was performed to create this manuscript.
\section{Conclusion}
AfterClick demonstrates, within its constructed asynchronous workflows, that local task success can hide forbidden effect history, enforcement delay can change outcomes with a fixed safety recommendation, serialized admission does not revoke previously admitted tickets, containment creates restoration obligations, and LLM review can be systematically overrestrictive relative to a declared rule. The most valuable insight is not that agents should reason faster or that a particular controller wins. It is that decisions, enforcement, immutable history and legitimate restoration require separate evidence. A bounded, traceable measurement study is justified by the current results; universal controller claims and broader causal explanations require further work.
\clearpage
\begin{thebibliography}{9}\small
\bibitem{undo} D. Sah, T. Sah, H. Jain, and T. Sah. UndoBench: Separating Task Competence from Recovery Capability in Tool-Using AI Agents. arXiv:2610.05622, 2026. \url{https://arxiv.org/abs/2610.05622}.
\bibitem{cordon} Z. Chen et al. Cordon: Semantic Transactions for Tool-Using LLM Agents. arXiv:2606.17573, 2026. \url{https://arxiv.org/abs/2606.17573}.
\bibitem{async} D. Zeng et al. Async-RBench: Benchmarking Asynchronous Task Execution in LLM Agents. arXiv:2610.32429, 2026. \url{https://arxiv.org/abs/2610.32429}.
\bibitem{rein} C. Duan et al. When to Intervene? An Empirical Study of Errors and Self-Recovery in LLM Agents. arXiv:2610.36471, 2026. \url{https://arxiv.org/abs/2610.36471}.
\bibitem{resume} G. Wu et al. Safe to Resume? Breaking Execution Continuity of Agent Execution via Rollback. arXiv:2608.29381, 2026. \url{https://arxiv.org/abs/2608.29381}.
\bibitem{agenttrust} C. Yang. AgentTrust: Runtime Safety Evaluation and Interception for AI Agent Tool Use. arXiv:2605.04785, 2026. \url{https://arxiv.org/abs/2605.04785}.
\bibitem{basin} D. Basin, V. Jug\'e, F. Klaedtke, and E. Z\u{a}linescu. Enforceable Security Policies Revisited. Principles of Security and Trust, 2012. \url{https://doi.org/10.1007/978-3-642-28641-4_17}.
\end{thebibliography}
\clearpage
\appendix
\section{Detailed experimental disposition and supplementary analyses}
\begin{table}[!htbp]\small\centering\caption{Complete original internal assignment accounting, by cohort and arm. Invalid and not-run cases retain unknown outcomes.}\label{tab:fullarms}
\begin{tabular}{@{}llrrrr@{}}\toprule
Cohort & Arm & Registered & Valid & Invalid & Not run\\\midrule
Primary & P0 clean & 432 & 141 & 28 & 263\\
Primary & P0 stress & 432 & 142 & 30 & 260\\
Primary & P1 stress & 432 & 142 & 30 & 260\\
Primary & P2 stress & 432 & 141 & 30 & 261\\
Primary & P3 stress & 432 & 136 & 34 & 262\\
Primary & P4 stress & 72 & 18 & 5 & 49\\
Primary & P5 stress & 72 & 24 & 4 & 44\\
Larger & P0 clean & 72 & 69 & 0 & 3\\
Larger & P0 stress & 72 & 65 & 1 & 6\\
Larger & P1 stress & 72 & 65 & 1 & 6\\
Larger & P2 stress & 72 & 65 & 2 & 5\\
Larger & P3 stress & 72 & 65 & 1 & 6\\
\bottomrule
\end{tabular}
\end{table}
\begin{table}[!htbp]\small\centering\caption{Matched policy contrasts in joint safe completion (pp). Source intervals condition on observed support. Degenerate intervals reflect the stratified sampling design.}\label{tab:contrasts}
\begin{tabular}{@{}llrrrr@{}}\toprule
Cohort & Contrast & Clusters & Estimate & 95\% CI & Assignment bounds\\\midrule
Primary & P3--P2 & 111 & $-13.51$ & $[-18.9,-9.0]$ & $[-72.5,63.4]$\\
Primary & P2--P1 & 112 & $2.68$ & $[-2.7,8.0]$ & $[-66.0,68.5]$\\
Primary & P1--P0 & 115 & $0.00$ & $[-5.7,5.2]$ & $[-67.1,67.1]$\\
Primary & P4--P1 & 16 & $31.25$ & $[31.2,31.2]$ & $[-63.9,83.3]$\\
Primary & P5--P1 & 20 & $-10.00$ & $[-10.0,-10.0]$ & $[-70.8,62.5]$\\
Primary & P5--P3 & 20 & $5.00$ & $[5.0,5.0]$ & $[-66.7,68.1]$\\
Larger & P3--P2 & 36 & $-6.94$ & $[-6.9,-6.9]$ & $[-15.3,4.2]$\\
Larger & P2--P1 & 36 & $5.56$ & $[5.6,5.6]$ & $[-5.6,13.9]$\\
Larger & P1--P0 & 36 & $-5.56$ & $[-5.6,-5.6]$ & $[-15.3,4.2]$\\
\bottomrule
\end{tabular}
\end{table}
The larger P4/P5 comparisons were not enrolled. Primary P3--P2 has 15 discordant task-mean signs, all negative, and the recorded source sign $p=0.000061$ conditional on observed support; larger P3--P2 has four discordant tasks (sign $p=0.125$). Neither conditional result repairs missing-assignment identification. Secondary tests retain original multiplicity correction. Mixed-model results with separation or optimization warnings are not used as headline evidence. The original frozen power analysis disclosed low power near zero success.
\fig{sensitivity}{.95}{Leave-one-application/family-out sensitivity of RQ1's observed conditional mismatch rate. Dashed lines give full observed cohort estimates. This is a structural sensitivity diagnostic, not a new independent sample.}
\subsection{External benchmark subset}
The official external cohort uses unchanged released tasks, tools, prompts and evaluator with five configured wrappers. It has 12 base tasks, two context seeds and 24 control plus 24 fault conditions per wrapper; all 240 assignments are terminal. Context seeds are not independently authored fault programs. Conditional recovery uses matched control-capable pairs, not all fault successes divided by total attempts. EvoUndo was unavailable; no fabricated value is assigned. Results do not establish general wrapper superiority or a replication of entire published systems.
\begin{table}[!htbp]\small\centering\caption{Official external subset: control and fault success, and conditional recovery on capable controls. 95\% intervals are exploratory task-cluster bootstrap intervals.}\label{tab:external}
\begin{tabular}{@{}lrrrrr@{}}\toprule
Wrapper & Control /24 & Fault /24 & Capable & Conditional & 95\% CI\\\midrule
Checkpoint & 20 & 6 & 20 & 30.0\% & 0.0--60.0\%\\
Idempotency & 19 & 8 & 19 & 42.1\% & 11.1--75.0\%\\
Naive & 20 & 5 & 20 & 25.0\% & 0.0--50.0\%\\
Saga & 21 & 6 & 21 & 28.6\% & 0.0--57.1\%\\
Vbr & 19 & 12 & 19 & 63.2\% & 35.0--87.5\%\\
\bottomrule
\end{tabular}
\end{table}
\subsection{Recoverability evidence and historical study scope}
\begin{table}[!htbp]\small\centering\caption{Selected prefix evidence labels. Categories shown are not exhaustive; no complex-prefix formal certificates exist.}\label{tab:prefix}
\begin{tabular}{@{}lrrrrr@{}}\toprule
Cohort & Nodes & Unknown & Safe initial witness & Safe final & Formal\\\midrule
Primary & 71138 & 59174 & 210 & 132 & 0\\
Larger & 36444 & 28512 & 57 & 35 & 0\\
\bottomrule
\end{tabular}
\end{table}
Earlier scripted and live studies include a seven-branch prototype, failed natural-task pilots and a 12-source prospective intervention study. In the latter, immediate protection has 11/12 joint successes and zero unauthorized effects, versus 10/12 and two effects for both delayed policies; the paired interval is wide and $p=1$. Earlier candidate controllers and negative results are retained as history rather than selected into a pooled effect estimate. Small finite-model certificates do not certify complex LLM prefixes. The present paper prioritizes the completed actuator/probe experiments and auditable V12 retrospective results.
\section{Evidence map and concrete episode provenance}\label{app:trace}
All results are grounded in repository revision \code{31eab61870befc520841d11f422bb1a372ba92b8}. Public source: \repo. The paper's plots and numerical tables are generated by direct reads from hash-verified compact reports at that revision; this presentation pipeline is not a raw-data reanalysis. The historical revision preserves all superseded source and paths; \path{archive/SNAPSHOTS.json} records the restoration assets. Detailed case cards, raw pointers, script hashes and the machine-generated claim registry remain the evidence trail, indexed by \path{evidence/index.json}.
\begin{table}[!htbp]\small\centering\caption{Core claims and compact-report paths in the current artifact.}\label{tab:evidence}
\begin{tabularx}{\linewidth}{@{}p{.28\linewidth}X@{}}\toprule
Claim & Repository path\\\midrule
Enrollment and validity & \path{research_ops/publication/PUBLICATION_STATUS.json} \\
State/history matrices, reviewer counts & \path{evidence/publication/analysis/report.json} \\
Controlled effect estimates and cells & \path{evidence/v12/controlled-actuator/report.json} \\
Actual restoration in scripted cohort & \path{evidence/publication/actuator-verification/report.json} \\
Historical late receipt witnesses & \path{evidence/publication/analysis/rq4-refined-ordering.csv} \\
Mechanism-probe registration and outcome & \path{evidence/publication/boundary/registration.json}; \path{evidence/publication/boundary/report.json} \\
Induced obligations & \path{evidence/publication/obligations/duties.csv}; \path{evidence/publication/checks/report.json} \\
Reviewer replay & \path{evidence/publication/replay/report.json} \\
Policy variants and uncertainty & \path{evidence/v12/primary-analysis/report.json}; \path{evidence/v12/larger-analysis/report.json} \\
Structural sensitivity & \path{evidence/publication/analysis/structural-sensitivity.csv} \\
External benchmark & \path{evidence/v12/official-external-analysis/report.json} \\
Claim mapping & \path{evidence/claim_evidence_matrix.csv}; \path{paper/claim_registry.json} \\\bottomrule
\end{tabularx}\end{table}
\textbf{Case A snapshot SHA-256:} \path{047b3db1b3760e85731a5ff5610d06264a2800e3aa13de839c042d469a536cef}. Case card \path{evidence/publication/analysis/cases/primary--TEST-ACCESS_MIGRATION-L1-FLOWBOARD-03-P0-tight-1-primary.md}. Pointers: \path{/v12_receipts/0}, \path{/v12_physical_effects/0}.

\textbf{Case B snapshot SHA-256:} \path{0f82e6a2dff1e93f2c6a3a47da2546482fe8fc66c01a946d88a7728812a4e31d}. Case card \path{evidence/publication/analysis/cases/larger--TEST-CUSTOMER_RELEASE-L1-OPSCONSOLE-05-P3-tight-1-larger.md}. Same receiver pointers; observed containment/restoration pointers \path{/v12_events/13} and \path{/v12_events/8} (pointer order is not assumed to be chronological index order).

\textbf{Forensic timeline:} raw final SHA-256 \path{da46276cc4bf43308efd542d4f91676add6fc1056da48f0c8455d11d1a4b7c29}. Browser-action SHA-256 \path{6ea2f9bde90330ea1491816d4cfa72af6d6a506a3551c6945be9fed9c1cd08d8}; commit-return pointer \path{/1/result/result/commit_upper_mono}. Reported timestamps in seconds: admission audit 252.349168596, protection audit 252.354895006, return upper bound 252.358651546, receipt 252.624891213. Figure~\ref{fig:forensic} subtracts the admission stamp and multiplies by 1000. This preserves the limits of the original clock interpretation.
\section{Reproduction protocol and research checklist}
The public repository includes compact reports, locked Python dependencies, browser clients, analysis programs and the current evidence registry. The original raw archive (637,853,954 bytes, 151,623 payload files) is locally retained; the audited 645,459,203-byte candidate adds audit metadata with original payloads unchanged. Original archive SHA-256: \path{c031afd3c46eab822770cc41e684d87f2a2aa786b64cff70eb79f90572667020}. Candidate SHA-256: \path{d82178a963b3fe1a4f39d72727c39a5648de8383e8e8bda09b8740779734d509}. The companion archive includes the publication-phase raw traces. No public download or DOI is claimed. The recorded candidate audit checks 263,779 expanded members and finds no credential-pattern matches; this is not a guarantee of privacy.

Use the current artifact and the separately supplied raw archives according to \path{docs/REPRODUCIBILITY.md}. Restore the exact pinned source with \path{archive/restore.py} for historical frozen analysis; the streamlined active tree is not the historical freeze. Compact verification and the repository's original paper generation require no model or provider calls:
\begin{verbatim}
uv sync --frozen --group research
uv run --frozen --group research python -m \
  research_ops.publication.program verify-data
uv run --frozen --group research python -m \
  research_ops.publication.program generate-paper
\end{verbatim}
Full analyses require the separately supplied archive and new output directories. With the original data restored as documented, run:
\begin{verbatim}
python -m research_ops.publication.program reproduce-existing \
  --historical-root /path/to/restored-history --output /tmp/ac-existing-new
python -m research_ops.publication.program analyze \
  --raw-root /path/to/restored-history --output /tmp/ac-analysis-new
\end{verbatim}
The live mechanism demonstration is distinct from saved-data analysis: start the isolated database, install Chromium, and invoke the documented \code{demo} entrypoint. No paid endpoint is needed. The historical \$18 budget and new \$28 ceiling are separate; the reported conservative totals are \$17.9673 and \$1.8895, respectively, with reported cash billing zero. These are guard accounting, not isolated invoice measurements. The old program's 41,732 model-call attempts and 309,714,800 known tokens include engineering, DEV and TEST; they are not completed episode counts. New replay accounts for 264 requests and 502,944 tokens; two setup calls are excluded from scientific replay enrollment.
\begin{table}[H]\small\centering\caption{Submission-readiness checklist (conference-style, not an official venue certification).}\label{tab:check}
\begin{tabularx}{\linewidth}{@{}p{.37\linewidth}X@{}}\toprule
Item & Status and implication\\\midrule
Research questions and outcome definitions & Specified; historical vs post hoc analyses distinguished.\\
Original enrollment and unknowns & Fully reconciled; original study incomplete.\\
Uncertainty and ablation limitations & Source-level intervals and unknown bounds reported; policy bundles not isolated.\\
Data and code access & Compact source public; raw release and licensing pending.\\
Human validation & Not performed; synthetic contracts and oracle agreement are automated.\\
Reproduction & Numerical Linux checks preserved; separate local historical cleanup regression passes 417/417.\\
Ethics and broader impacts & Synthetic effects only; privacy review of release still required.\\
Compute and failed runs & Separate ledgers, negative runs and unknowns retained.\\
Anonymity and venue policy & Public repository identifies the owner; adapt artifact for double-blind review if required.\\\bottomrule
\end{tabularx}\end{table}
The supplied draft relied on compact evidence. Subsequent repository cleanup reverified raw-data analyses and ran the complete historical regression, documented separately in \path{archive/audit/VALIDATION.json}; independent human review remains absent. All measurements are inherited from the pinned repository's compact evidence and its documented analyses. New figures and interpretive hypotheses do not represent new experiments. This draft uses a conference-inspired single-column layout, not an official NeurIPS submission template or page-limit certification.
\end{document}
